\begin{question}{27}{
    De Landmacht onderzoekt de effectiviteit van een nieuw soort camouflagenetten die Fennek-verkenningsvoertuigen moeten beschermen tegen detectie door moderne sensoren en drones.
    Deze camouflagenetten worden in drie terreintypes (bebost gebied, duinlandschap en stedelijk gebied) getest op basis van in totaal $300$ verkenningsvluchten door een vijand-simulerende drone-eenheid.
    Per vlucht wordt geregistreerd of het voertuig wordt gedetecteerd of niet, zoals te zien in onderstaande tabel:
    \begin{center}
        \renewcommand{\arraystretch}{1.25}
        \begin{tabular}{cccc}
            \toprule
                & \textbf{Detectie} & \textbf{Geen detectie} & \textbf{Totaal} \\
            \midrule
                Bebost gebied & $26$ & $84$ & $110$  \\
                Duinlandschap & $48$ & $72$ & $120$  \\
                Stedelijk gebied & $32$ & $38$ & $70$  \\
            \midrule
                \textbf{Totaal} & $106$ & $194$ & $300$ \\
            \bottomrule
        \end{tabular}
    \end{center}

    De Landmacht wil onderzoeken of een Fennek-voertuig in een bepaald terreintype eerder wordt gedetecteerd.
}

    \subquestion{3}{
        Welke hypothesetoets zou hiervoor geschikt zijn?
        Formuleer $H_0$ en $H_1$ voor deze toets.
    }

    \solution{
        Om te bepalen of er een statistisch significant verschil is in detectiekansen tussen de drie terreintypen, kunnen we een chi-kwadraattoets voor onafhankelijkheid gebruiken. \rubric{2}
        De nulhypothese $H_0$ en de alternatieve hypothesetoets $H_1$ kunnen dan als volgt worden geformuleerd:
        \begin{description}
            \item[$H_0$:] de detectiekans is niet afhankelijk van het terreintype. \rubric{1}
            \item[$H_1$:] de detectiekans is wel afhankelijk van het terreintype. \rubric{1}
        \end{description}
    }

    \subquestion{6}{
        Bereken de toetsingsgrootheid van deze hypothesetoets.
    }
    \solution{
        Om de verwachte frequenties $E_{ij}$ te bepalen in een chi-kwadraattoets voor onafhankelijkheid, moeten we gebruik maken van de volgende formule:
        \[
            E_{ij} = \frac{(\text{rijtotaal}_i) \cdot (\text{kolomtotaal}_j)}{\text{totaal}}
        \]
        We berekenen de verwachte frequenties voor elke cel in de tabel:
        \begin{center}
            \renewcommand{\arraystretch}{1.25}
            \begin{tabular}{cccc}
                \toprule
                    & \textbf{Detectie} & \textbf{Geen detectie} & \textbf{Totaal} \\
                \midrule
                    Bebost gebied & $38.8667$ & $71.1333$ & $110$ \\ 
                    Duinlandschap & $42.4$ & $77.6$ & $120$ \\ 
                    Stedelijk gebied & $24.7333$ & $45.2667$ & $70$ \\
                \midrule
                    \textbf{Totaal} & $106$ & $194$ & $300$ \rubric{3} \\  
                \bottomrule 
            \end{tabular}
        \end{center}

        Vervolgens berekenen we de toetsingsgrootheid $X^2$ als volgt:

        \begin{align*}
            \chi^2 &= \frac{(O_{1, 1} - E_{1,1})^2}{E_{1, 1}} + \frac{(O_{1, 2} - E_{1,2})^2}{E_{1, 2}} + \ldots + \frac{(O_{3, 2} - E_{3,2})^2}{E_{3, 2}} \\
                    &= \frac{(26 - 38.8667)^2}{38.8667} + \frac{(84 - 71.1333)^2}{71.1333} + \ldots + \frac{(38 - 45.2667)^2}{45.2667} \\
                    &\approx 11.032. \rubric{3}
        \end{align*}
    
        
    }

    \subquestion{5}{
        Wat is de kansverdeling van de toetsingsgrootheid onder $H_0$? Bepaal met behulp van deze kansverdeling de $p$-waarde.
    }
    \solution{
        De toetsingsgrootheid $X^2$ volgt onder de nulhypothese een $\chi^2$-verdeling met $\text{df}= (\#\text{rijen}-1)\cdot(\#\text{kolommen}-1) = (3-1)\cdot(2-1) = 2$ vrijheidsgraden. \rubric{2}
        De $p$-waarde is de kans op een nog extremere uitkomst voor de toetsingsgrootheid, in het geval van de chi-kwadraatverdeling een nog grotere uitkomst. \rubric{1}
        De $p$-waarde behorende bij onze toetsingsgrootheid is daarom gelijk aan        
        \begin{align*}
            p = P(X^2 > 11.032) &= \chi^2\text{cdf}(\text{lower}=11.032, \text{upper}=10^{99}, \text{df}=2) \\
                                &\approx 0.004.   \rubric{2}
        \end{align*}
    }

    \subquestion{5}{
        Formuleer een conclusie voor deze toets (op basis van de berekende $p$-waarde en een significantieniveau $\alpha = 0.05$) en vertaal deze terug naar de originele probleemcontext.
        Geef een verklaring voor deze conclusie aan de hand van bovenstaande tabel.
    }
    \solution{
        Aangezien de $p$-waarde kleiner is dan het significantieniveau $\alpha = 0.05$, wordt de nulhypothese $H_0$ verworpen. \rubric{1}
        Er is voldoende reden om aan te nemen dat de detectiekans afhankelijk is van het terreintype waarin de Fennek-voertuigen zich bevinden.\rubric{2}

        Als we kijken naar de geobserveerde en verwachte frequenties, dan zien we dat het percentage detecties in stedelijke gebied en het duinlandschap hoger is dan verwacht, terwijl het percentage detecties in bebost gebieden lager is dan verwacht. \rubric{2}
    }
    
    \subquestion{8}{
        We bekijken nu de detectiekans $p$ over alle terreintypes gezamenlijk.
        De fabrikant van de camouflagenetten beweert dat de detectiekans \SI{30}{\percent} is.
        Bereken een \SI{95}{\percent}-betrouwbaarheidsinterval voor $p$ met behulp van de Clopper-Pearson methode.
        Wat kunnen we concluderen over de claim van de fabrikant?
    }

    \solution{
        Als we over alle terreintypes gezamenlijk kijken, zien we dat $106$ van de $300$ verkenningsvluchten (steekproef!) heeft geresulteerd in een detectie. \rubric{1}
        Op basis van deze aantallen gaan we een betrouwbaarheidsinterval bepalen voor de algemene detectiekans $p$ (populatieparameter!).

        Omdat de gewenste betrouwbaarheid \SI{95}{\percent} is, geldt dat $\alpha = 0.05$.

        De Clopper-Pearson methode werkt als volgt:
        \begin{enumerate}
            \item Bepaal de succeskans $p_1$ waarvoor geldt dat de linkeroverschrijdingskans van de uitkomst $k=106$ gelijk is aan $\alpha/2$, oftewel $P(X \leq 106) = \alpha/2 = 0.025$.
            Voer hiervoor in het solver menu van de grafische rekenmachine in:
            \begin{align*}
                E_1 &= \binomcdf(n=300; p_1=X; k=106) \\
                E_2 &= 0.025 \rubric{1}
            \end{align*}
            De solver optie geeft een waarde van $p_1 \approx 0.4103$.\rubric{1}
            \item Bepaal de succeskans $p_2$ waarvoor geldt dat de rechteroverschrijdingskans van de uitkomst $k=106$ gelijk is aan $\alpha/2$, oftewel $P(X \geq 106) = 1 - P(X \le 105) = \alpha/2 = 0.025$.
            Voer hiervoor in het solver menu van de grafische rekenmachine:
            \begin{align*}
                E_1 &= 1 - \binomcdf(n=300, p_1=X; k=105) \\
                E_2 &= 0.025 \rubric{1}
            \end{align*}
            De solver optie geeft een waarde van $p_2 \approx 0.2993$.\rubric{1}
        \end{enumerate}
        We vinden het Clopper-Pearson interval door de twee gevonden waarden als grenzen te nemen, oftewel het \SI{95}{\percent}-betrouwbaarheidsinterval voor de detectiekans $p$ van Fennek-voertuigen met de nieuwe camouflagenetten is
        is gelijk aan $[0.2993, 0.4103]$. \rubric{1}

        We zien dat de geclaimde succeskans $p = 0.3$ wel in dit interval ligt. \rubric{1}
        Er is dus onvoldoende reden om aan te nemen dat de succeskans significant hoger is dan \SI{30}{\percent}. \rubric{1}
    }
\end{question}